\documentclass{article}

\textwidth=6.5in
\textheight=8.5in
\voffset=-54pt
\oddsidemargin=0in
\evensidemargin=0in
\setlength{\parskip}{8pt}
\setlength{\parindent}{0pt}

\usepackage{workbook}

\usepackage[sols]{handout}

\begin{document}

\begin{center}
  \large\textsc{Problem Set 32 - Area and Density Problems}
  
      \pspace{12pt}
  \begin{problemonly}\large Name: \rule{5cm}{0.15mm}\\ \end{problemonly}
\end{center}

\begin{grading}

    \begin{enumerate}
    \item 5 points
    \item 10 points
    \item 15 points
    \item 10 points
    \item 3 points
    \end{enumerate}
  Total: 43 points
\end{grading}

\begin{enumerate}



  \item

\begin{grading}
    5 points. 2 points for finding the necessary intersection points, and 3 points for
writing and evaluating the integrals. This problem does not need to be done with slicing.
\end{grading}

\begin{problem}[\vfill]
    The picture below shows the graphs of $y = \sqrt{3} \cos x$ and $y = -
    \sin x$.  Find the shaded area.

    \emph{Hint:} As you can see, there are two pieces to account for;
    you'll need to write separate integrals for the two pieces (do you see
    why?).  \emph{Do} evaluate your integrals to get a numerical answer.

    \begin{tikzpicture}
      \begin{axis}[stack plots=y, xtick={-3.14, 3.14, 6.28}, ytick={-1,
          1}, xticklabels={$-\pi$, $\pi$, $2\pi$}, disabledatascaling,
        xmin=-3.2, width=3in, height=2in]  
        \addplot+[domain=-pi/3:2*pi/3] {-sin(deg(x))}; %
        \addplot+[domain=-pi/3:2*pi/3, fillregion]
        {sqrt(3)*cos(deg(x))+sin(deg(x))} \closedcycle; %
        \resetstackedplots %
        \addplot+[domain=2*pi/3:5*pi/3] {sqrt(3)*cos(deg(x))}; %
        \addplot+[domain=2*pi/3:5*pi/3, fillregion]
        {-sqrt(3)*cos(deg(x))-sin(deg(x))} \closedcycle; %
        \resetstackedplots %
        \addplot+[domain=-pi:2*pi+.2] {-sin(deg(x))}; %
        \resetstackedplots %
        \addplot+[domain=-pi:2*pi+.2] {sqrt(3)*cos(deg(x))}; %
      \end{axis}
    \end{tikzpicture}
  \end{problem}

  \begin{solution}
    First, let's find where the two curves intersect on $[-\pi, 2\pi]$.
    We need to solve $\sqrt{3} \cos x = - \sin x$.  Dividing both sides by
    $-\cos x$, we can rewrite this as $-\sqrt{3} = \tan x$, which has
    solutions $x = - \frac{\pi}{3}$, $\frac{2 \pi}{3}$, and $\frac{5
      \pi}{3}$ on $[-\pi, 2\pi]$.  Let's label these in our picture:

    \begin{tikzpicture}
      \begin{axis}[stack plots=y, xtick={-1.0472, 2.0944, 5.23599},
        xticklabels={$-\frac{\pi }{3}$, $\frac{2 \pi }{3}$, $\frac{5 \pi
          }{3}$}, 
        ytick=\empty,
        xmin=-3.2, width=3in, height=2in, clip=false]  
        \addplot+[domain=-pi/3:2*pi/3, draw=none] {-sin(deg(x))}; %
        \addplot+[domain=-pi/3:2*pi/3, fillregion, draw=none]
        {sqrt(3)*cos(deg(x))+sin(deg(x))} \closedcycle; %
        \resetstackedplots %
        \addplot+[domain=2*pi/3:5*pi/3, draw=none] {sqrt(3)*cos(deg(x))}; %
        \addplot+[domain=2*pi/3:5*pi/3, fillregion, draw=none]
        {-sqrt(3)*cos(deg(x))-sin(deg(x))} \closedcycle; %
        \resetstackedplots %
        \addplot+[domain=-pi:2*pi+.2, red] {-sin(deg(x))} node[right]{$y =
          - \sin x$}; %
        \resetstackedplots %
        \addplot+[domain=-pi:2*pi+.2, blue] {sqrt(3)*cos(deg(x))}
        node[right]{$y = \sqrt{3} \cos x$}; %
      \end{axis}
    \end{tikzpicture}

    We need to write two separate integrals because the two regions have
    different upper functions and different lower functions.  On $\left[ -
      \frac{\pi}{3}, \frac{2 \pi}{3} \right]$, the upper function is
    $\sqrt{3} \cos x$ and the lower function is $-\sin x$, so the area of
    that region is $\displaystyle \int_{-\pi/3}^{2 \pi/3} \left( \sqrt{3}
      \cos x + \sin x \right) dx$.  On $\left[ \frac{2 \pi}{3}, \frac{5
        \pi}{3} \right]$, the upper function is $-\sin x$, and the lower
    function is $\sqrt{3} \cos x$, so the area of that region is
    $\displaystyle \int_{2 \pi/3}^{5 \pi/3} \left( -\sin x - \sqrt{3} \cos
      x\right) dx$.   So, the area we want is
    \begin{align*}
      & \int_{-\pi/3}^{2 \pi/3} \left( \sqrt{3} \cos x + \sin x \right) dx
      + \int_{2 \pi/3}^{5 \pi/3} \left( -\sin x - \sqrt{3} \cos x\right) dx
      \\
      &= \left( \sqrt{3} \sin x - \cos x \Big|_{-\pi/3}^{2\pi/3} \right) +
      \left( \cos x - \sqrt{3} \sin x \Big|_{2 \pi/3}^{5 \pi/3} \right)
      \\
      &= \left( \left[ \sqrt{3} \left( \frac{\sqrt{3}}{2} \right) +
          \frac{1}{2} \right] - \left[ \sqrt{3} \left( -
            \frac{\sqrt{3}}{2} \right) - \frac{1}{2} \right] \right) +
      \left( \left[ \frac{1}{2} - \sqrt{3} \left( - \frac{\sqrt{3}}{2}
          \right) \right] - \left[ -\frac{1}{2} - \sqrt{3} \left(
            \frac{\sqrt{3}}{2} \right) \right] \right) \\
      &= \boxed{8}
    \end{align*}
  \end{solution}

\item 

\begin{problem}[\nstretch{1.5}]
    Find the area of the region bounded by $y = \ln x$, $x =
    e^2$, and $y = 0$.

    \emph{Hint:} You'll have to decide whether to slice vertically or
    horizontally.  The two methods will give you different integrals, and
    only one of those integrals is easy to evaluate.  So, if the first way
    you slice gives you a difficult integral, try slicing the other way!
  \end{problem}

  \begin{solution}
    Here's a picture of the region: 
    \begin{center}
      \begin{tikzpicture}
        \begin{axis}[xmin=0, xtick={1, 7.38906}, xticklabels={$1$, $e^2$},
          ymin=-1, width=2.5in, ymax=2.25] 
          \addplot+[domain=1:e^2, fillregion] {ln(x)} \closedcycle;
          \addplot+[domain=0.1: 8] {ln(x)};
        \end{axis}
      \end{tikzpicture}
    \end{center}
    If we were to slice vertically, the integral we'd end up with is
    $\displaystyle \int_1^{e^2} \ln x\,dx$, which would require integration by parts.  Instead, let's
    slice horizontally as this makes the resulting computations more straightforward.

    The top point of the region is $\left( e^2, \ln \left( e^2 \right)
    \right) = \left( e^2, 2 \right)$, so we need to slice $[0, 2]$ on the
    $y$-axis into $n$ horizontal slices of equal height $\Delta y$.  We
    label $y_0 = 0, \ldots, y_n = 2$ as usual; let's focus on the $k$-th
    slice.  We can approximate this slice as a rectangle:
    \begin{center}
      \begin{tikzpicture}
        \begin{axis}[axis on top=true, xmin=0, xtick={1, 7.38906},
          xticklabels={1, $e^2$}, ytick={0, 0.2857, ..., 2.01},
          yticklabels={$y_0 = 0$,,,,$y_k$,,,$y_n = 2$}, width=2.5in,
          ymin=-1, axis y line=left, ymax=2.25, tick label
          style={fill=none}] 
          \addplot+[domain=1:e^2, fillregion, draw=none] {2}
          \closedcycle; %
          \addplot+[xcomb] coordinates { (e^2, 0) (e^2, 1*2/7) (e^2,
            2*2/7) (e^2, 3*2/7) (e^2, 4*2/7) (e^2, 5*2/7) (e^2, 6*2/7)
            (e^2, 2) }; %
          \addplot+[fill=white, draw=none, domain=1:e^2] {ln(x)} -- (0,
          7.38906) -- (0, 0) -- (1, 0); %
          \addplot+[domain=0.5:8] {ln(x)}; %
          \draw[dashed] (0, 8/7) -- ({e^(8/7)}, 8/7); %
          \draw[fill=cyan, draw=blue] ({e^(4*2/7)}, 4*2/7) -- (e^2, 4*2/7)
          -- (e^2, 3*2/7) -- ({e^(4*2/7)}, 3*2/7) -- cycle; %
          \draw (e^2, 0) -- (e^2, 2);
          \addplot+[domain=0.1: 8] {ln(x)};
        \end{axis}
      \end{tikzpicture}
    \end{center}
    This rectangle has height $\Delta y$.  Its right side is on the line
    $x = e^2$, and its left side is on the curve $y = \ln x$ (which can be
    rewritten as $x = e^y$); therefore, its width is (right
    $x$-coordinate) $-$ (left $x$-coordinate) $= e^2 - e^{y_k}$.  So,
    \[ \text{area of $k$-th slice} \approx \left( e^2 - e^{y_k} \right)
    \Delta y.\]  Summing over all slices and taking the limit as the
    number of slices $\to \infty$, the area of the region is
    \begin{align*}
      \int_0^2 \left( e^2 - e^y \right) dy &= e^2 y - e^y \Big|_0^2 \\
      &= \left( 2e^2 - e^2 \right) - (0 - 1) \\
      &= \boxed{e^2 + 1}
    \end{align*}
  \end{solution}

\item

\begin{grading}
   15 points: 3/3/2/2/5
\end{grading}

For help with this problem, read pp. 827-835 in Gottlieb. Note the list of common
  errors on page 835!

  A plot on a farm is rectangular with dimensions of
  $300$ meters by $200$ meters.  A straight irrigation pipe runs down one
  300-meter side of the plot.  The crop yield density decreases as the
  distance from the irrigation pipe increases.  Suppose that the yield
  density is given by $\rho(x)$ kg/m$^2$, where $x$ is the distance in
  meters from the irrigation pipe.  We'd like to find the total yield (in
  kg) of the plot.

  \begin{enumerate}

  \item
    \begin{problem}[\stretch{2}]
      Show in a sketch the meaning of $x$.  (That is, draw in an $x$-axis
      and tell us the value of $x$ at a few places; for example, where is
      $x = 0$?)  Then, show how you can \textbf{slice} the plot to solve
      this problem.
    \end{problem}

    \begin{solution}
      Since $x$ measures distance from the pipe, $x = 0$ at the pipe and
      $x = 200$ at the opposite end of the field.  Here's one possible
      sketch, with the pipe in blue:
      \begin{center}        
        \begin{tikzpicture}
          \begin{axis}[xmin=0, xmax=220, ymin=0, ymax=310, hide y axis,
            xtick={0, 200}, axis equal=true, axis x line=bottom, 
            x=0.015cm, y=0.015cm, xlabel=$x$]
            \draw (0, 10) -- (200, 10) -- (200, 310) -- (0, 310) -- cycle; 
            \draw[ultra thick, blue] (0, 10) -- (0, 310);
          \end{axis}
        \end{tikzpicture}
      \end{center}
      Since density varies with distance to the pipe, we should slice
      parallel to the pipe so that the density within each slice is
      approximately constant:
      \begin{center}        
        \begin{tikzpicture}
          \begin{axis}[xmin=0, xmax=220, ymin=0, ymax=310, hide y axis,
            xtick={0, 20, ..., 200}, xticklabels={0,,,,,,,,,,200}, axis
            equal=true, axis x line=bottom,  
            x=0.015cm, y=0.015cm, xlabel=$x$]
            \pgfplotsforeachungrouped \x in {20, 40, ..., 200}
            {
              \edef\temp{
                \noexpand\draw (\x, 10) -- (\x, 310);
              }
              \temp
            }            
            \draw (0, 10) -- (200, 10) -- (200, 310) -- (0, 310) -- cycle; 
            \draw[ultra thick, blue] (0, 10) -- (0, 310);
          \end{axis}
        \end{tikzpicture}
      \end{center}
      (Note that we're slicing the interval $[0, 200]$.)
    \end{solution}

  \item
    \begin{problem}[\vfill]
      Write an expression that gives the \textbf{approximate} crop yield
      in the $k$-th slice. 
    \end{problem}

    \begin{solution}
      We're focusing on the $k$-th slice:
      \begin{center}        
        \begin{tikzpicture}
          \begin{axis}[xmin=0, xmax=220, ymin=0, ymax=310, hide y axis,
            xtick={0, 20, ..., 200},
            xticklabels={$x_0$,,,,,,,$x_k$,,,$x_n$}, axis equal=true, axis
            x line=bottom, x=0.015cm, y=0.015cm, xlabel=$x$]
            \draw[fill=cyan] (120, 10) -- (120, 310) -- (140, 310) --
            (140, 10) -- cycle;
            \pgfplotsforeachungrouped \x in {20, 40, ..., 200}
            {
              \edef\temp{
                \noexpand\draw (\x, 10) -- (\x, 310);
              }
              \temp
            }            
            \draw (0, 10) -- (200, 10) -- (200, 310) -- (0, 310) -- cycle; 
            \draw[ultra thick, blue] (0, 10) -- (0, 310);
          \end{axis}
        \end{tikzpicture}
      \end{center}
      The crop yield in this slice can be approximated as
      \begin{align*}
        \text{yield of $k$-th slice} & \approx (\text{yield density in
          $k$-th slice}) \cdot (\text{area of $k$-th slice}) \\
        \intertext{The slice is a rectangle of width $\Delta x$ and height
          $300$, so}
        & \approx \boxed{\rho(x_k) \cdot 300 \Delta x}
      \end{align*}
    \end{solution}

  \item
    \begin{problem}[\vfill]
      Write a general Riemann \textbf{sum} (that is, a Riemann sum with
      $n$ terms) that approximates the total crop yield from the plot.
    \end{problem}

    \begin{solution}
      To approximate the yield of the whole plot, we just add up the
      yields of all the slices; thus, the yield of the whole plot is 
      $\approx \boxed{\sum_{k=1}^n 300 \rho(x_k) \Delta x}$. 
    \end{solution}

  \item
    \begin{problem}[\vfill]
      Write an integral that gives the exact total yield from the plot.
      (Your integral will be in terms of the density function $\rho(x)$.) 
    \end{problem}

    \begin{solution}
      Taking the limit as $n \to \infty$ of our Riemann sum, the exact
      yield from the plot is $\boxed{ \int_0^{200} 300 \rho(x)\,dx}$.
    \end{solution}

\pspace{\newpage}

  \item
    \begin{problem}[\vfill]
      If $\displaystyle \rho(x) = \frac{1000}{(x + 100)^2}$, evaluate your
      integral.  (You should get 2000 kg.) 
    \end{problem}

    \begin{solution}
      We're evaluating
      \begin{align*}
        \int_0^{200} 300 \cdot \frac{1000}{(x + 100)^2}\,dx 
        &= 300,000 \int_0^{200} \frac{1}{(x + 100)^2}\,dx \\
        \intertext{Let's substitute $u = x + 100$; then, $du = dx$.  As
          $x$ goes from $0$ to $200$, $u = x + 100$ goes from $0 + 100 =
          100$ to $200 + 100 = 300$:} %
        &= 300,000 \int_{100}^{300} \frac{1}{u^2}\,du \\
        &= 300,000 \int_{100}^{300} u^{-2}\,du \\
        &= 300,000 \left( - u^{-1} \right) \Big|_{100}^{300} \\
        &= 300,000 \left( - \frac{1}{300} + \frac{1}{100} \right) \\
        &= 2000,
      \end{align*}
      so the field yields 2000 kg of corn.
    \end{solution}

  \end{enumerate}

\item

\begin{grading}
    10 points: 5 per part
\end{grading}

A consultant suggests instead running a straight irrigation pipe $300$ meters
  long down the center of the plot, dividing it in half lengthwise,
  would increase yield by 50\%. Show that this is true. 
  Suppose the crop yield density decreases as the distance from the 
  irrigation pipe increases and that the yield density is given by 
  $\rho(x)$ kg/m$^2$, where $x$ is the distance in meters from the
  irrigation pipe.

  \begin{enumerate}

  \item
    \begin{problem}[\stretch{0.5}]
      Write an integral that gives the total yield from the plot.
    \end{problem}

    \begin{solution}
      Here's a sketch.  Notice that, since $x$ represents distance from
      the pipe, $x$ increases from 0 to 100 when you go left \textit{or}
      right from the pipe.
      \begin{center}
        \includegraphics[width=1.5in]{images/density-rectangle-2}
      \end{center}
      Because of symmetry, the yield in the left half of the plot is the
      same as the yield in the right half.  Therefore, we only need to
      figure out the yield in one half; let's do the right half.

      Since the yield depends on the distance from the irrigation pipe, we
      want to make slices in which the distance to the irrigation pipe is
      approximately constant.  Therefore, we'll slice parallel to the
      pipe:
      \begin{center}
        \includegraphics[width=1.5in]{images/density-rectangle-2-sliced}
      \end{center}
      This amounts to slicing the interval $[0, 100]$ into $n$ pieces, and
      we label $x_0 = 0, \ldots, x_n = 100$ as usual.

      The area of each slice is $300 \Delta x$ (the width is $\Delta x$,
      the height is 300).  In the $k$-th slice, the yield density is
      approximately $\rho(x_k)$ grams per square meter, so the yield is
      approximately $\rho(x_k) \cdot 300 \Delta x = 300 \rho(x_k) \Delta
      x$.  Summing over all $n$ slices gives us the Riemann sum
      $\displaystyle \sum_{k=1}^n 300 \rho(x_k) \Delta x$ as an
      approximation for the yield of half of the plot.  Therefore, the
      total yield is approximated by \fbox{$\displaystyle 2 \sum_{k=1}^n
        300 \rho(x_k) \Delta x$}.

      Taking the limit of the Riemann sum as $n \to \infty$, the yield for
      the plot is \fbox{$\displaystyle 2 \int_0^{100} 300 \rho(x)\,dx$}.
    \end{solution}

  \item
    \begin{problem}[\vfill\newpage]
      If $\displaystyle \rho(x) = \frac{1000}{(x + 100)^2}$, evaluate your
      integral to get the total yield from the plot.  (You should get 3000
      kg.)
    \end{problem}

    \begin{solution}
      We need to evaluate
      \begin{align*}
        2 \int_0^{100} 300 \rho(x)\,dx &= 600,000 \int_0^{100} \frac{1}{(x
          + 100)^2}\,dx \\
        \intertext{Let's substitute $u = x + 100$.  Then, $du = dx$; as
          $x$ goes from $0$ to $100$, $u = x + 100$ goes from $0 + 100 =
          100$ to $100 + 100 = 200$.}
        &= 600,000 \int_{100}^{200} u^{-2}\,du \\
        &= -600,000 u^{-1} \Big|_{100}^{200} \\
        &= 3000,
      \end{align*}
      so the total yield from this plot is $3000$ kg. 
    \end{solution}

  \end{enumerate}

% \item ({Problem Set 30}, {{\#3}})

%  A farmer has a plot of land which is 100 meters $\times$ 100 meters.
%   He installs an irrigation pipe running along the southern border of the
%   land (shown in blue below).  He plans to plant several different types
%   of crops on this plot of land; one of the crops will be kale, which he
%   decides to plant in a 50 meter $\times$ 25 meter rectangle at the
%   northwest corner of the plot, with the 50-meter sides parallel to the
%   pipe (the shaded area below).

%   \begin{tikzpicture}[x=1cm, y=1cm]
%     \draw[fill=white!90!black, draw=gray] (0, 4) -- (2, 4) -- (2, 3)
%     -- (0, 3) -- (0, 4); %
%     \draw[blue, ultra thick] (0, 0) -- (4, 0); %
%     \draw (4, 0) -- (4, 4) -- (0, 4) -- (0, 0); %
%   \end{tikzpicture}    

%   How well the kale grows depends on how far it is from the irrigation
%   pipe; suppose that the farmer can grow $\rho(x) = \dfrac{80}{4x + 25}$
%   kg of kale per square meter, where $x$ is the distance in meters from
%   the irrigation pipe.  We'd like to figure out how much kale the farmer
%   will grow.

%   \begin{enumerate}
%   \item
%     \begin{problem}
%       Show in a sketch the meaning of $x$.  
%     \end{problem}

%     \begin{solution}
%       Since $x$ measures distance from the pipe, $x = 0$ at the pipe and
%       $x = 100$ on the opposite end of the plot.  Therefore, the side of
%       the kale patch closer to the pipe is at $x = 75$:
%       \begin{center}
%         \begin{tikzpicture}
%           \begin{axis}[hide x axis, axis y line=left, ylabel=$x$,
%             ytick={0, 3, 4}, yticklabels={0, 75, 100}, xmin=0, xmax=4,
%             ymin=0, ymax=4.25, clip=false, x=1cm, y=1cm]
%             \begin{scope}[xshift=8pt]
%               \draw[fill=white!90!black, draw=gray] (0, 4) -- (2, 4) --
%               (2, 3) -- (0, 3) -- (0, 4); %
%               \draw[blue, ultra thick] (0, 0) -- (4, 0); %
%               \draw (4, 0) -- (4, 4) -- (0, 4) -- (0, 0); %
%             \end{scope}
%           \end{axis}
%         \end{tikzpicture}
%       \end{center}
%     \end{solution}

%   \item
%     \begin{problem}
%       Find the amount of kale that the farmer will grow.  (You should find
%       that he grows about $268$ kg of kale, but please give an exact
%       answer.) 
%     \end{problem}

%     \begin{solution}
%       Since the density of kale depends on the distance from the pipe, we
%       must slice parallel to the pipe in order to have the density of kale
%       within each slice be approximately constant.  So, we'll slice the
%       interval $[75, 100]$ on the $x$-axis into $n$ pieces of equal
%       thickness $\Delta x$.  We'll label the endpoints $x_0 = 0, \ldots,
%       x_n = 30$ and focus on the $k$-th slice; here's a picture, showing
%       only the kale patch:
%       \begin{center}
%         \begin{tikzpicture}
%           \begin{axis}[hide x axis, axis y line=left, ylabel=$x$,
%             ytick={3, 3.14157, ..., 4.01}, yticklabels={$x_0 =
%               75$,,,,,$x_k$,,$x_n = 100$}, xmin=0, xmax=4,
%             ymin=3, ymax=4.25, clip=false, x=2cm, y=2cm]
%             \begin{scope}[xshift=8pt]
%               \draw[fill=yellow, draw=none] (0, {25/7}) -- (2, {25/7}) --
%               (2, {26/7}) -- (0, {26/7}) -- cycle;
%               \draw[fill=white!90!black, draw=gray] (0, 4) -- (2, 4) --
%               (2, 3) -- (0, 3) -- (0, 4); %
%               \draw[fill=yellow, draw=none] (0, {25/7}) -- (2, {25/7}) --
%               (2, {26/7}) -- (0, {26/7}) -- cycle; %
%               \draw (0, 4) -- (2, 4) -- (2, 3) -- node[below] {50} (0, 3)
%               -- (0, 4); %
%               \pgfplotsforeachungrouped \x in {21, 22, ..., 28} {
%                 \edef\temp{ \noexpand\draw (0, {\x/7}) -- (2, {\x/7}); }
%                 \temp }
%             \end{scope}
%           \end{axis}
%         \end{tikzpicture}
%       \end{center}
%       Then,
%       \begin{align*}
%         \text{amount of kale in $k$-th slice} & \approx (\text{density of
%           kale in $k$-th slice})(\text{area of $k$-th slice}) \\
%         & \approx \rho(x_k) \cdot 50 \Delta x
%       \end{align*}
%       Summing over all slices and taking the limit as the number of slices
%       $\to \infty$, the amount of kale the farmer will grow is
%       \begin{align*}
%         \int_{75}^{100} \rho(x) \cdot 50\,dx &= 4000 \int_{75}^{100}
%         \frac{1}{4x + 25}\,dx \\
%         \intertext{You can substitute $u = 4x + 25$ to do this integral,
%           or you might be able to just play around to see that an
%           antiderivative of $\dfrac{1}{4x + 25}$ is $\dfrac{1}{4} \ln |4x
%           + 25|$:}
%         &= \left. 4000 \cdot \frac{1}{4} \ln |4x + 25| \right|_{75}^{100}
%         \\
%         &= \boxed{1000 \ln 425 - 1000 \ln 325 \text{ kg}} \\
%         \intertext{We can use log rules to rewrite this a little:}
%         &= 1000 \ln \left( \frac{425}{325} \right) \\
%         &= \boxed{1000 \ln \left( \frac{17}{13} \right) \text{ kg}}
%       \end{align*}
%     \end{solution}

%   \end{enumerate}

%   \emph{Thought question: How would your integral change if the
%       farmer instead decided to plant the kale in the southeast corner of
%       his plot (still with the 50-meter sides parallel to the pipe)?}

% \item ({Problem Set 30}, {{\#4}})

%   A plot of land in a rain forest is shaped like a right triangle with
%   sides of length $1$, $\sqrt{3}$, and 2 miles.  The $\sqrt{3}$-mile side
%   lies along a river, and vegetation density varies with $x$, distance to
%   the river.  Suppose vegetation density is given by $\rho(x)$ million
%   plants per square mile.

%   \begin{enumerate}
%   \item
%     \begin{problem}
%       Write a general Riemann sum that approximates the amount of
%       vegetation in the plot.
%     \end{problem}

%     \begin{solution}
%       Here's a picture of the plot, with the river on the left side (in
%       blue).  Since $x$ represents distance to the river, $x = 0$ at the
%       river, and $x = 1$ at the vertex furthest from the river.
%       \begin{center}
%         \begin{tikzpicture}
%           \begin{axis}[xmin=0, xmax=1.2, ymin=0, ymax=1.8, clip=false,
%             hide y axis, axis x line=bottom, x=2cm, y=2cm, xlabel=$x$,
%             xtick={0, 1}]
%             \draw[yshift=8pt] (0, 0) -- node[left] {$\sqrt{3}$} (0,
%             {sqrt(3)}) -- (1, 0) -- (0, 0);
%             \draw[yshift=8pt, ultra thick, blue] (0, 0) -- (0, {sqrt(3)});
%           \end{axis}
%         \end{tikzpicture}
%       \end{center}
%       Since the vegetation density varies with distance to the river, we
%       should slice parallel to the river in order to get slices where the
%       density within any given slice is approximately constant.  So, let's
%       slice $[0, 1]$ into $n$ pieces of equal thickness $\Delta x$ and
%       focus on the $k$-th slice:
%       \begin{center}
%         \begin{tikzpicture}[declare function={f(\x)=sqrt(3)*(1-\x);}]
%           \begin{axis}[xmin=0, xmax=1.2, ymin=0, ymax=1.8, clip=false,
%             hide y axis, axis x line=bottom, x=2cm, y=2cm, xlabel=$x$,
%             xtick={0, 0.125, ..., 1.01},
%             xticklabels={$x_0$,,,,$x_k$,,,,$x_n$}]
%             \begin{scope}[yshift=8pt]
%               \draw (0, 0) -- node[left] {$\sqrt{3}$} (0, {sqrt(3)}) --
%               (1, 0) -- (0, 0); %
%               \draw[fill=cyan] (0.375, 0) -- (0.5, 0) -- (0.5, {f(0.5)})
%               -- (0.375, {f(0.375)}) -- cycle; %
%               \draw[ultra thick, blue] (0, 0) -- (0, {sqrt(3)}); %
%               \pgfplotsforeachungrouped \x in {0, 0.125, ..., 1.01} {%
%                 \edef\temp{%
%                   \noexpand\draw (\x, 0) -- (\x, {sqrt(3)*(1-\x)}); %
%                 }%
%                 \temp%
%               }
%             \end{scope}
%           \end{axis}
%         \end{tikzpicture}
%       \end{center}
%       We can approximate the amount of vegetation in the $k$-th slice as
%       \[ \text{amount of vegetation in $k$-th slice} \approx
%       (\text{density of vegetation in $k$-th slice})(\text{area of $k$-th
%         slice}).\] The density of vegetation in the $k$-th slice is
%       approximately $\rho(x_k)$.  To deal with the area, we can
%       approximate the slice as a rectangle of width $\Delta x$, but we
%       still need to express its (approximate) height in terms of $x_k$.
%       We can use similar triangles to do that.  Let $h_k$ be the
%       approximate height of the $k$-th slice:
%       \begin{center}
%         \begin{tikzpicture}[declare function={f(\x)=sqrt(3)*(1-\x);}]
%           \begin{axis}[xmin=0, xmax=1.2, ymin=0, ymax=1.8, clip=false,
%             hide x axis, hide y axis, axis x line=bottom, x=2cm, y=2cm,
%             xlabel=$x$, xtick={0, 0.125, ..., 1.01},
%             xticklabels={$x_0$,,,,$x_k$,,,,$x_n$}]
%             \begin{scope}[yshift=8pt]
%               \draw (0, 0) -- node[left] {$\sqrt{3}$} (0, {sqrt(3)}) --
%               (1, 0) -- (0, 0); %
%               \draw[fill=cyan] (0.375, 0) -- (0.5, 0) -- node[right, inner
%               sep=1pt]{$h_k$} (0.5, {f(0.5)}) -- (0.375, {f(0.375)}) --
%               cycle; %
%               \draw[ultra thick, blue] (0, 0) -- (0, {sqrt(3)}); %
%               \draw[decorate, decoration={brace, mirror}, yshift=-2pt]
%               (0.5, 0) -- node[below, yshift=-2pt] {$1 - x_k$} (1, 0); %
%               \draw[yshift=-24pt, |-|] (0, 0) -- node[fill=white] {1} (1,
%               0); 
%             \end{scope}
%           \end{axis}
%         \end{tikzpicture}
%       \end{center}
%       By similar triangles, $\displaystyle \frac{\sqrt{3}}{1} =
%       \frac{h_k}{1 - x_k}$, so $h_k = \sqrt{3} (1 - x_k)$.  Therefore,
%       \begin{align*}
%         \text{amount of vegetation in the $k$-th slice} & \approx
%         (\text{density of vegetation in $k$-th slice})(\text{area of $k$-th
%           slice}) \\
%         &\approx \rho(x_k) \cdot (\text{height of $k$-th
%           slice})(\text{width of $k$-th slice}) \\
%         & \approx \rho(x_k) \cdot \sqrt{3} (1 - x_k) \Delta x
%       \end{align*}
%       Summing over all slices, the amount of vegetation in the plot is
%       approximately $\boxed{\sum_{k=1}^n \sqrt{3} (1 - x_k) \rho(x_k)
%         \Delta x}$.     
%     \end{solution}

%   \item
%     \begin{problem}
%       Find a definite integral that gives the exact amount of vegetation
%       in the plot.
%     \end{problem}

%     \begin{solution}
%       Taking the limit as $n \to \infty$ of our Riemann sum, the amount of
%       vegetation in the plot is exactly $\boxed{ \int_0^1 \sqrt{3} (1 - x)
%         \rho(x)\,dx}$.
%     \end{solution}

%   \item
%     \begin{problem}
%       If $\rho(x) = 3 - x$ milion plants per square mile, evaluate your
%       integral to find the exact amount of vegetation in the plot.  (You
%       should get $\frac{4 \sqrt{3}}{3}$ million plants.) 
%     \end{problem}

%     \begin{solution}
%       We need to evaluate
%       \begin{align*}
%         \int_0^1 \sqrt{3} (1 - x)(3-x)\,dx &= \sqrt3 \int_0^1
%         (x^2-4x+3)\,dx \\
%         &= \left. \sqrt3\left(\frac{1}{3}x^3 - 2x^2+3x\right)\right|_0^1 \\
%         &= \boxed{\frac{4 \sqrt{3}}{3}}
%       \end{align*}
%     \end{solution}

%   \end{enumerate}

% \item ({Problem Set 30}, {{\#5}})

%   ($\star$) A farmer has a square plot of land with area 2 square
%   kilometers; each side of the square is $\sqrt{2}$ km long.  The farmer
%   is deciding how to position an irrigation pipe on this square plot; he
%   is considering two options.  One option is to position the irrigation
%   pipe so that it runs from the center of one side of the square plot to
%   the center of the opposite side of the square plot. The other option is
%   to position the irrigation pipe so that it runs diagonally from one
%   corner of the square plot to the opposite corner of the square
%   plot. (See the illustrations below.)  In both cases, the density of crop
%   yield is given by $\rho(s)= 100-30 s^2$ kg/km$^2$, where $s$ is the
%   distance to the irrigation pipe.

%   \emph{Note:} Try to take advantage of symmetry in this problem!

%   \begin{center}
%     \begin{tabular}{c c c}
%       \begin{tikzpicture}[x=2cm, y=2cm]
%         \draw[color=white] (0, {sqrt(2)/2-1}) -- (0, {sqrt(2)/2+1});
%         \draw (0, 0) -- (0, {sqrt(2)}) -- node[midway, yshift=0.5em]{$\sqrt{2}$} ({sqrt(2)}, {sqrt(2)}) --
%         ({sqrt(2)}, 0) -- (0, 0);
%         \draw[dashed] (0, {sqrt(2)/2}) -- node[midway, yshift=0.5em]{pipe}
%         ({sqrt(2)}, {sqrt(2)/2});
%       \end{tikzpicture}
%       &
%       \hspace{0.25in} & 
%       \begin{tikzpicture}[x=2cm, y=2cm]
%         \draw[color=white] (0, {sqrt(2)/2-1}) -- (0, {sqrt(2)/2+1}); %
%         \draw (0, 0) -- (0, {sqrt(2)}) -- node[midway,
%         yshift=0.5em]{$\sqrt{2}$} ({sqrt(2)}, {sqrt(2)}) -- ({sqrt(2)}, 0)
%         -- (0, 0); %
%         \draw[dashed] (0, {sqrt(2)}) -- ({sqrt(2)}, {0}); %
%         \draw[decorate, decoration={text along path, text={|\small|pipe}, text
%           align={center}, raise=4pt}] (0, {sqrt(2)}) -- ({sqrt(2)},
%         {0}); %
%       \end{tikzpicture}
%       \\
%       Option 1 & & Option 2 
%     \end{tabular}
%   \end{center}

%   \begin{enumerate}
%   \item
%     \begin{problem}
%       Find the crop yield if the farmer uses the first option. 
%     \end{problem}

%     \begin{solution} % originally by Jameel, modified by Janet
%       By symmetry, the yield is twice the yield for the top half of the
%       field.  The left picture below is a picture of the top half,
%       together with an ``$s$-axis'' ($s$ measures distance from the pipe):
%       \begin{center}
%         \includegraphics[height=1in]{images/fall2008-midterm1-2-option1-axis}
%         \hspace{0.5in}
%         \includegraphics[height=1in]{images/fall2008-midterm1-2-option1-slice} 
%       \end{center}
%       We will slice horizontally so that the crop density on each slice is
%       approximately constant.  This means that we'll slice the interval
%       $\left[0, \frac{\sqrt{2}}{2} \right]$ into $n$ pieces of equal
%       height $\Delta s$, and we'll label $s_0 = 0, \ldots, s_n =
%       \frac{\sqrt{2}}{2}$ as usual.  Then, we focus on the $k$-th slice,
%       shown in the right picture above.
%       \begin{align*}
%         \text{yield of the $k$-th slice} & \approx \text{(approximate
%           yield density in $k$-th slice)} \cdot
%         \text{(area of $k$-th slice)} \\
%         & \approx \rho(s_k) \sqrt{2} \Delta s
%       \end{align*}
%       Summing over all slices and taking the limit as the number of slices
%       $\to \infty$, the yield in the top half of the field is
%       $\displaystyle \int_0^{\sqrt{2} / 2} \sqrt{2} \rho(s)\,ds$.
%       Therefore, the yield in the entire field is twice that, which is
%       \begin{align*}
%         2\int_0^{\sqrt{2}/2} (100-30 s^2) \sqrt{2} \,ds &=
%         2\sqrt{2}\int_0^{\sqrt{2}/2} (100 - 30s^2) ds \\
%         & = 2 \sqrt{2} \left( 100s - 10 s^3 \right) \Big|_0^{\sqrt{2} / 2} \\ 
%         & = \boxed{190 \text{ kg}}
%       \end{align*}
%     \end{solution}

%   \item
%     \begin{problem}
%       Find the crop yield if the farmer uses the second option. 
%     \end{problem}

%     \begin{solution} % by Jameel, modified by Janet
%       The pipe divides the plot into two halves; by symmetry, each half
%       has the same yield, so we can look at just one half.  Since $s$
%       measures distance to the pipe, $s$ is 0 along the pipe and increases
%       as we move away from the pipe.  So, here's one half of the plot
%       (we've rotated the picture just to make it easier to look at):
%       \begin{center}
%         \begin{tikzpicture}
%           \begin{axis}[x=2cm, y=2cm, hide y axis, xlabel=$s$, xmin=0,
%             xmax=1.2, ymin=0, ymax=2, clip=false, xtick=\empty]
%             \draw[dashed] (0, 1) -- node[left] {pipe} (0, -1); %
%             \draw (0, 1) -- node[anchor=south west] {$\sqrt{2}$} (1, 0) --
%             (0, -1);
%           \end{axis}
%         \end{tikzpicture}
%       \end{center}
%       By the Pythagorean Theorem, the pipe's length is $2$.  In addition,
%       since the yield density depends only on distance from the pipe, the
%       top half of this piece has the same yield as the bottom half;
%       therefore, we can look just at the top half (which is one quarter of
%       the whole field).  Here's a picture of that, together with an 
%       $s$-axis:
%       \begin{center}
%         \begin{tikzpicture}
%           \begin{axis}[x=2cm, y=2cm, clip=false, xmin=0, xmax=1.2, ymin=1,
%             ymax=2, hide y axis, xlabel=$s$, axis x line=bottom, xtick={0,
%               1}] 
%             \begin{scope}[yshift=8pt]
%               \draw[dashed] (0, 2) -- node[left] {pipe} (0, 1); %
%               \draw (0, 2) -- node[anchor=south west] {$\sqrt{2}$} (1, 1)
%               -- (0, 1); %
%             \end{scope}
%           \end{axis}
%         \end{tikzpicture}
%       \end{center}
%       We want to slice so that the density in each slice is approximately
%       constant. Since the density depends only on distance from the pipe,
%       we should slice parallel to the pipe.  This means we'll be slicing
%       the interval $[0, 1]$ (on the $s$-axis) into $n$ pieces; we'll label
%       $s_0 = 0, \ldots, s_n = 1$ as usual and focus on the $k$-th slice:
%       \begin{center}
%         \begin{tikzpicture}
%           \begin{axis}[xmin=0, xmax=1.2, ymin=0, ymax=1.2, clip=false,
%             hide y axis, axis x line=bottom, x=2cm, y=2cm, xlabel=$s$,
%             xtick={0, 0.125, ..., 1.01},
%             xticklabels={$s_0$,,,,,$s_k$,,,$s_n$}] 
%             \begin{scope}[yshift=8pt]
%               \draw[draw=none, fill=cyan] (0.5, 0) -- (0.5, 0.5) --
%               (0.625, 0.375) -- (0.625, 0) -- cycle; %
%               \draw (0, 1) -- (1, 0) -- (0, 0); %
%               \pgfplotsforeachungrouped \x in {0.125, 0.25, ..., 1.01} {%
%                 \edef\temp{%
%                   \noexpand\draw (\x, 0) -- (\x, {-1/1*(\x - 1)}); %
%                 }%
%                 \temp%
%               } \draw[dashed] (0, 0) -- node[left] {1} (0, 1);
%             \end{scope}
%             \begin{scope}[font=\footnotesize]
%               \draw (0, -13pt) node[rotate=90] {$=$};
%               \draw (1, -13pt) node[rotate=90] {$=$};
%               \draw (0, -21pt) node {$0$};
%               \draw (1, -21pt) node {$1$};
%             \end{scope}
%           \end{axis}
%         \end{tikzpicture}
%       \end{center}
%       The yield in the $k$-th slice can be approximated as $(\text{yield
%         density in $k$-th slice})(\text{area of $k$-th slice})$.  The
%       yield density in the $k$-th slice is approximately $\rho(s_k)$.  To
%       deal with the area, we can approximate the slice as a rectangle of
%       width $\Delta s$, but we still need to express its (approximate)
%       height in terms of $s_k$.  We can use similar triangles to do that.
%       Let $h_k$ be the approximate height of the $k$-th slice:
%       \begin{center}
%         \begin{tikzpicture}
%           \begin{axis}[xmin=0, xmax=1.2, ymin=0, ymax=1.2, clip=false,
%             hide y axis, axis x line=bottom, x=2cm, y=2cm,
%             hide x axis] 
%             \begin{scope}[yshift=8pt]
%               \draw[fill=cyan] (0.5, 0) -- (0.5, 0.5) -- (0.625, 0.375) --
%               node[right, black, inner sep=1pt] {$h_k$} (0.625, 0) --
%               cycle; %
%               \draw (0, 1) -- (1, 0) -- (0, 0); %
%               \draw[dashed] (0, 0) -- node[left] {1} (0, 1); %
%               \draw[decorate, decoration={brace, mirror}, yshift=-2pt]
%               (0.625, 0) -- node[below, yshift=-2pt] {$1 -
%                 s_k$} (1, 0); \draw[decorate, decoration={brace, mirror},
%               yshift=-24pt] (0, 0) -- node[below, yshift=-2pt] {$1$} (1, 0);
%             \end{scope}
%           \end{axis}
%         \end{tikzpicture}
%       \end{center}
%       By simliar triangles, $\displaystyle \frac{h_k}{1 - s_k} =
%       \frac{1}{1}$, so $h_k = 1 - s_k$.  Therefore,
%       \begin{align*}
%         \text{yield in $k$-th slice} & \approx (\text{yield density in
%           $k$-th slice})(\text{area of $k$-th
%           slice}) \\
%         & \approx \rho(s_k) h_k \Delta s \\
%         & \approx \rho(s_k) \left( 1 - s_k \right) \Delta s
%       \end{align*}
%       Summing over all slices and taking the limit as the number of slices
%       $\to \infty$, the yield for one quarter of the plot is
%       $\displaystyle \int_0^{1} \rho(s) (1 - s)\,ds$.  Therefore, the
%       total yield is
%       \begin{align*}
%         \boxed{4 \int_0^{1} \rho(s) (1 - s)\,ds} &= 4 \int_0^1 (100 -
%         30s^2)(1 - s)\,ds \\
%         &= 4 \int_0^1 ( 100 - 100 s - 30s^2 + 30s^3 )\,ds \\
%         &= 4 \left( \left. 100 s - 50 s^2 - 10 s^3 + \frac{15}{2} s^4
%           \right|_0^1 \right) \\
%         &= \boxed{190 \text{ kg}}
%       \end{align*}
%     \end{solution}

%   \end{enumerate}

%   \begin{problemonly}
%     If you've done everything correctly, you should find that the two
%     options give the same yield.  (This is a coincidence, not something
%     that is always true.)
%   \end{problemonly}

% Based on from Robin's Spring 2009 1b file Review_exam_1.tex, #1;
%   apparently written by Jeechul Woo

\item \emph{Which way to slice?}
Suppose you have a flat platter shaped like a disk of radius 25 cm, and
  the density of cobalt in the glaze of the platter varies. 
  (Cobalt creates a blue color.) In each part below, you're
  given a different function $\rho$ describing the density of cobalt
  in the
  platter, in mg/cm$^2$.  In each case, decide whether you would slice
  into concentric rings or into strips in order to approximate the amount
  of cobalt in the platter. (Drawing a picture and labeling the distance $x$ may help.)

\probonly{
  \begin{center}
    \begin{tabular}{c c c}
      \begin{tikzpicture}
        \foreach \x in {0.2, 0.4, ..., 1.61}
        {
          \draw (0, 0) circle (\x);
        }
      \end{tikzpicture} & 
      \hspace{0.25in}
      &
      \begin{tikzpicture}
        \draw (0, 0) circle (1.6);
        \foreach \x in {-1.4, -1.2, ..., 1.41}
        {
          \draw (\x, {-sqrt(1.6^2-abs(\x)^2)}) -- (\x, {sqrt(1.6^2-abs(\x)^2)});
        }
      \end{tikzpicture}
      \\
      Concentric rings & & Strips 
    \end{tabular}
  \end{center}
}

  \begin{enumerate}
  \item 
    \begin{problem}[\vfill]
      $\rho(x) = x^2 + 1$, where $x$ is the distance to the center point
      of the platter.
    \end{problem}

    \begin{solution}
      We need to slice into \fbox{concentric rings} in order for the
      density within any slice to be approximately constant, since the distance $x$ is measured radially outwards from the center.
      
    \begin{tikzpicture}
        \fill (0,0) circle (2pt);
        \draw[very thick,->] (0,0) -- (0:1.8) node[pos=1.1] {$x$};
        \draw[very thick,->] (0,0) -- (120:1.8) node[pos=1.1] {$x$};
        \draw[very thick,->] (0,0) -- (240:1.8) node[pos=1.1] {$x$};
        \foreach \x in {0.2, 0.4, ..., 1.61}
        {
          \draw (0, 0) circle (\x);
        }
      \end{tikzpicture}
    \end{solution}

  \item
    \begin{problem}[\vfill]
      $\rho(x) = x^2 + 1$, where $x$ is the distance to the vertical diameter of the platter.
    \end{problem}

    \begin{solution}
      We need to slice into vertical \fbox{strips} in order
      for the density within any slice to be approximately constant, since the distance $x$ is measured horizontally outwards from the vertical diameter.

          \begin{tikzpicture}
        \draw[very thick] (0,1.6) -- (0,-1.6);
        \draw[very thick,->] (0,0) -- (1.8,0) node[pos=1.1] {$x$};
        \draw[very thick,->] (0,0) -- (-1.8,0) node[pos=1.1] {$x$};
        \draw (0, 0) circle (1.6);
        \foreach \x in {-1.4, -1.2, ..., 1.41}
        {
          \draw (\x, {-sqrt(1.6^2-abs(\x)^2)}) -- (\x, {sqrt(1.6^2-abs(\x)^2)});
        }
      \end{tikzpicture}
    \end{solution}

  \item
    \begin{problem}[\vfill]
      $\rho(x) = x^2 + 1$, where $x$ is distance to the edge of the
      platter. 
    \end{problem}

    \begin{solution}
      We need to slice into \fbox{concentric rings} in order for the
      density within any slice to be approximately constant, since the distance $x$ is measured radially inwards from the edge of the platter.
      
          \begin{tikzpicture}
        \draw[very thick] (0,0) circle (1.6);
        \draw[very thick,->] (0:1.6) -- (0:0) node[above, pos=0.9] {$x$};
        \draw[very thick,->] (120:1.6) -- (120:0) node[above right, pos=0.9] {};
        \draw[very thick,->] (240:1.6) -- (240:0) node[below right, pos=0.9] {};
        \foreach \x in {0.2, 0.4, ..., 1.61}
        {
          \draw (0, 0) circle (\x);
        }
      \end{tikzpicture}
    \end{solution}

  \end{enumerate}

\pspace{\newpage}

\item 
\begin{problem}
    Next class, we will construct and evaluate integrals based on slicing a circular region into thin rings. The mathematical name for a circular ring-shaped region is an \emph{annulus}. To prepare, let's investigate the area of an annulus like the one shown in the diagram. Since it's a small area, we'll call it $\Delta A$.
\end{problem}

\begin{center}
\hfill
    \begin{tikzpicture}[scale=0.75]

  % Parameters
  \def\rin{1.5}    % Inner radius
  \def\rout{2}     % Outer radius

  % Draw outer circle
  \filldraw[draw=black, thick, fill=gray!20] (0,0) circle (\rout);
  % Cut inner circle to create annulus
  \filldraw[draw=black, thick, fill=white] (0,0) circle (\rin);

  % Draw center
  \fill (0,0) circle (2pt);

% Set angles for each radius line
  \def\anglein{0}
  \def\angleout{90}

  % Draw inner radius line
  \draw[thick, densely dashed] (0,0) -- ({\anglein}:\rin)
    node[midway, below] {\large$r_1$};

  % Draw outer radius line
  \draw[thick, densely dashed] (0,0) -- ({\angleout}:\rout)
    node[midway, right] {\large$r_2$};

\end{tikzpicture}
\hfill
    \begin{tikzpicture}[scale=0.75]

  % Parameters
  \def\rin{1.5}    % Inner radius
  \def\rout{2}     % Outer radius
  \def\ravg{1.75}  % Average radius
  \def\drangle{0} % Angle for Δr label
    \def\labeldist{0.8} % Additional radial distance for the Δr label

  % Draw outer circle
  \filldraw[draw=black, thick, fill=gray!20] (0,0) circle (\rout);
  % Cut inner circle to create annulus
  \filldraw[draw=black, thick, fill=white] (0,0) circle (\rin);
  % Draw average radius circle
  \draw[black, thick, densely dashed] (0,0) circle (\ravg);
    % Draw center
  \fill (0,0) circle (2pt);
  
% Set angles for each radius line
  \def\anglein{0}
  \def\angleout{90}

  % Draw width indicator (thickness of annulus)
  \draw[ultra thick] ({\drangle}:\rin) -- ({\drangle}:\rout);

  % Add external label and pointer arrow for Δr
  \node[align=center] (deltar) at (20:\rout + \labeldist) {\large$\Delta r$};
  \draw[->, thick] (deltar) -- ({\drangle}:\ravg);

  % Draw average radius
  \draw[thick, densely dashed] (0,0) -- (0,\ravg)
    node[midway, right] {\large$r$};

\end{tikzpicture}
\hfill\strut
\end{center}

\begin{enumerate}
    \item 
    \begin{problem}
        Find the area of the annulus $\Delta A$ in terms of the inner radius $r_1$ and the inner radius $r_2$.
    \end{problem}
    \pspace{\vfill}

    \begin{solution}
    The area of the annulus is the area of a circle of radius $r_2$ minus the area of a circle of radius $r_1$. So $\Delta A = \pi r_2^2 - \pi r_1^2$.
    \end{solution}

    \item
    \begin{problem}
        Since we're going to use these areas in our Riemann sums, it will be useful for us to express $\Delta A$ in terms of $\Delta r = r_2 - r_1$. Algebraically manipulate your formula from part (a) to show that $\Delta A = 2\pi r \Delta r$, where $r$ is the average radius, $r=\frac{1}{2}(r_1+r_2)$.
    \end{problem}
    \pspace{\stretch{2}}

    \begin{solution}
        \begin{align*}
            \Delta A &= \pi r_2^2 - \pi r_1^2 \\
            &= \pi (r_2^2 - r_2^2) \\
            &= \pi (r_2+r_1)(r_2-r_1) \\
            &= 2\pi \cdot \tfrac{1}{2}(r_2+r_2)\cdot (r_2-r_1) \\
            &= 2\pi r\Delta r
        \end{align*}
    \end{solution}

        \item
    \begin{problem}
        Informally, we can imagine ``cutting and unrolling'' the annulus into a ``rectangle.'' Explain why this matches the area formula $\Delta A = 2\pi r\Delta r$. Drawing a diagram may help.
    \end{problem}
    \pspace{\stretch{2}}

    \begin{solution}
    From the diagram below, we can see that if we ``cut and unroll'' the annulus into a ``rectangle,'' then the height of the rectangle is $\Delta r$ and the length is $2\pi r$, giving an area of $\Delta A = 2\pi r \Delta r$. Even though the annulus cannot actually be unrolled into a rectangle, this gives us a heuristic to remember the area formula.
    
        \begin{tikzpicture}
  % Rectangle dimensions
  \def\length{11}
  \def\width{0.5}

  % Draw rectangle
  \filldraw[fill=gray!20, draw=black, thick] (0,0) rectangle (\length,\width);

  % Label length (bottom center)
  \draw[<->] (0,-0.3) -- (\length,-0.3)
    node[midway, below] {\large$2\pi r$};

  % Label width (right side)
  \draw[<->] (\length + 0.3, 0) -- (\length + 0.3, \width)
    node[midway, right] {\large$\Delta r$};

\end{tikzpicture}
    \end{solution}
\end{enumerate}

\end{enumerate}

\spaceonly{\psreflection{
Optional reading for next class is \S 31.1 in Gottlieb.
}}

\end{document}